% !TeX spellcheck = en_US
\documentclass[aspectratio=1610]{beamer}
\mode<presentation>
{
  \usetheme{Boadilla}
  \setbeamercovered{transparent}
  
  \definecolor{cobalt}{rgb}{0.0, 0.20, 0.50}
  
  \usecolortheme[named=cobalt]{structure}
  %\usecolortheme{lily}
  %\setbeamercolor{headline}{fg=cyan!80!black}
  % \setbeamercolor{frametitle}{bg=white}.
}

\makeatletter
\define@key{beamerframe}{wide}[10pt]{%
	\def\beamer@cramped{\itemsep #1\topsep0.5pt\relax}}
\makeatother
	

\usepackage[english]{babel}
\usepackage[utf8]{inputenc}
\usepackage{times}
\usepackage[T1]{fontenc}
\usepackage{array}
\usepackage{xcolor}
\usepackage{colortbl}
\usepackage{booktabs}

\title[Winter - Lecture 8] % (optional, use only with long paper titles)
{Macroeconomic Theory and Policy: Lecture  W8 (Ch.14)}

\author[University of Toronto] % (optional, use only with lots of authors)
{}
\date[] % (optional, should be abbreviation of conference name)

\begin{document}

\begin{frame}
  \titlepage
\end{frame}

\begin{frame}[wide]{What are You Going to Learn (Ch. 14)}
	\begin{itemize}
		\item Various policy responses to the financial crisis
		\item The Fed’s balance sheet
		\item Contributors to a very large negative shock in AD
		\begin{itemize}
			\item Gap between the fed funds rate and market interest rates
			\item A household balance  sheet crisis
			\item Substantial uncertainty
		\end{itemize}
		\item Unconventional actions of the Fed
		
	\end{itemize}
\end{frame}

\begin{frame}[wide]{Policy Responses to the Financial Crisis}
	\begin{itemize}
		\item Simple policy rule: the central bank rate as a function of the gap between the current inflation rate and target rate
		\item The Taylor rule expands upon this: allows the current level of short-run output to influence the fed funds rate
		\begin{itemize}
			\item e.g., $i_t= 1\% + 1.5 \pi_t + 0.5 \widetilde{Y_t}$
		\end{itemize}
	\end{itemize}
\end{frame}

\begin{frame}[wide]{The Fed Funds Rate and the Taylor Rule}
	\begin{itemize}
		\item Some people argue that leading up to the crisis, the fed funds rate may have been too low for too long
		\item But, during the crisis, the actual interest rate that firms faced was high
		\begin{itemize}
			\item from 2002 to 2006 the fed funds rate was lower than the Taylor rule suggested
		\end{itemize}
		\item After the crisis, Taylor rule suggests a higher interest rate than the Fed Rate
	\end{itemize}
	\begin{figure}
		\centering
		\includegraphics[width=0.7\linewidth]{"Figures/fig 14_04"}
	\end{figure}
	
\end{frame}

\begin{frame}[wide]{The Fed Funds Rate and the Taylor Rule}
	\begin{itemize}
		\item Output was below potential throughout 2002 and into 2003
		\begin{itemize}
			\item led to inflation being well below its target, and contributed to the federal funds rate remaining low
		\end{itemize}
		\item Short run output has since grown beyond potential, and contributed to the increase in the fed funds rate
	\end{itemize}
	\begin{figure}
		\centering
		\includegraphics[width=0.7\linewidth]{"Figures/fig 14_05"}
	\end{figure}
\end{frame}

\begin{frame}[wide]{The Fed Funds Rate and the Taylor Rule}
	\begin{itemize}
		\item Decline in inflation in 2002 up to the end of 2003
		\item Fed was concerned about the possibility of deflation at that time
		\item The decline of inflation	was not nearly as sharp as a standard Phillips curve would have predicted during the crisis 
	\end{itemize}
	\begin{figure}
		\centering
		\includegraphics[width=0.7\linewidth]{"Figures/fig 14_06"}
	\end{figure}
	
\end{frame}

\begin{frame}[wide]{How Large is the Output Gap?}
	\begin{itemize}
		\item The closing of the gap happened because of revised potential GDP, not from a return to the previous GDP trend
		\item Why has potential GDP fallen? Low investment rates and TFP growth slowed
	\end{itemize}
	\begin{figure}
		\centering
		\includegraphics[width=0.5\linewidth]{"Figures/fig 14_07"}
	\end{figure}
\end{frame}

\begin{frame}[wide]{The Money Supply}
	\begin{itemize}
		\item Friedman and Schwartz claim that excessively tight monetary policy by the Federal Reserve and the ensuing deflation was the principal cause of the Great Depression
		\vspace{1mm}
		\begin{itemize}
			\item They cite that the money supply declined sharply between 1929 and 1933
		\end{itemize}
		\item Fed is currently focused on stimulating the economy and preventing deflation
		\vspace{1mm}
		\begin{itemize}
			\item Expanded the money supply at the end of 2008
		\end{itemize}
	\end{itemize}
\end{frame}

\begin{frame}[wide]{The Great Depression in the IS/MP Framework}
	\begin{itemize}
		\item 
	\end{itemize}
	\begin{figure}
		\centering
		\includegraphics[width=0.7\linewidth]{"Figures/fig 14_13"}
	\end{figure}
	
\end{frame}

\begin{frame}[wide]{The Phillips Curve and the Great Depression}
	\begin{itemize}
		\item 
	\end{itemize}
	\begin{figure}
		\centering
		\includegraphics[width=0.7\linewidth]{"Figures/fig 14_14"}
	\end{figure}
	
\end{frame}

\begin{frame}[wide]{The Growth Rate of Various Money Supply Measures}
	\begin{itemize}
		\item Currency is the narrowest measure, followed by M1 and M2
		\item All exhibit rapid growth by the end of 2008
		\item From the end of 2008 to 2009, the Fed was focused on increasing the money supply, stimulating the economy, and avoiding deflation
	\end{itemize}
	\begin{figure}
		\centering
		\includegraphics[width=0.5\linewidth]{"Figures/fig 14_08"}
	\end{figure}
	
\end{frame}

\begin{frame}[wide]{The Fed’s Balance Sheet}
	\begin{itemize}
		\item When conventional monetary policy failed, central bank adopt a set of unconventional monetary policy
		\begin{itemize}
			\item Allowed these institutions to swap less liquid financial instruments for Treasury securities
			\item The provision of liquidity to the large government-sponsored mortgage companies Fannie Mae and Freddie Mac
			\item Direct capital injections through the Troubled Asset Relief Program (TARP) 
		\end{itemize} 
	\end{itemize}
\end{frame}

\begin{frame}[wide]{The Fed’s Balance Sheet}
	\begin{itemize}
		\item These unconventional monetary policy dramatically change the Fed’s balance sheet
		\begin{itemize}
			\item Most assets were U.S. Treasury securities while most liabilities were currency held by the public
		\end{itemize}
		\item The composition of assets and liabilities also changed
		\begin{itemize}
			\item On the asset side: the Fed expanded its assets to include lending (through loans or securities) to the rest of the economy, not only to financial institutions but also to nonfinancial corporations
			\item On the liability side: additional lending was not financed by printing money.
			The Fed borrowed from the Treasury and bank excess reserves
		\end{itemize}
	\end{itemize}
\end{frame}

\begin{frame}[wide]{The Federal Reserve’s Balance Sheet (billions of dollars)}
	\begin{itemize}
		\item Various programs are combined into Mortgage-Backed Securities and Other
		\item The increase are mostly financed by an increase in commercial bank's reserve
		\begin{itemize}
			\item Fed bought commercial paper or securitized loans from financial institutions and paid for these purchases by crediting their reserve accounts
		\end{itemize}
		\item But banks kept reserves with the Fed instead of loaning them out
		\begin{itemize}
			\item Beside that bank were reluctant to make risky loans,the Fed began paying interest on excess reserves -- Central bank can adjust this interest rate and control the money flow
		\end{itemize}
	\end{itemize}
	\begin{figure}
		\centering
		\includegraphics[width=0.7\linewidth]{"Figures/fig 14_09"}
	\end{figure}
	
\end{frame}

\begin{frame}[wide]{Financial Assets of the Federal Reserve}
	\begin{itemize}
		\item Purchased a range of securities and made a large volume of loans 
		\item Quantitative easing: central bank purchases of assets other than short-term government bonds
		\item Initially, the Fed acts as the ``lender of last resort''. Then, bought only MBS in QE1 and bought long-term U.S. Treasuries in QE2. During QE3, they bought both
	\end{itemize}
	\begin{figure}
		\centering
		\includegraphics[width=0.5\linewidth]{"Figures/fig 14_12"}
	\end{figure}
\end{frame}

\begin{frame}[wide]{Should Monetary Policy Respond to Asset Prices?}
	\begin{itemize}
		\item It is now evident that there was a housing bubble in the mid-2000s
		\item Should the Fed attempt to curtail asset price bubbles when they occur? The Fed Chairman at that time Bernanke argued in 2000
		\begin{itemize}
			\item Bubbles are difficult to identify in real time
			\item Monetary policy is too imprecise to deal with bubbles
			\item Policymakers should use more refined instruments:
			\begin{itemize}
				\item capital requirements
				\item regulation of lending standards
			\end{itemize}
		\end{itemize}
	\end{itemize}
\end{frame}

\begin{frame}[wide]{The Price–Earnings Ratio in the Stock Market}
	\begin{itemize}
		\item Often serve as a measure of bubbles
		\item In asset pricing models, stock prices should be equal to the present discounted value of future profits
		\item Whenever the ratio of stock prices to company earnings gets too far away from its mean, it tends to revert back
	\end{itemize}
	\begin{figure}
		\centering
		\includegraphics[width=0.5\linewidth]{"Figures/fig 14_10"}
	\end{figure}
\end{frame}

\begin{frame}[wide]{Real Home Prices in the United States}
	\begin{itemize}
		\item By nearly any measure, the appreciation of home prices between 2000 and 2006 was unusual
		\item Moreover, previous episodes of rapidly rising real home prices were often followed by real declines
	\end{itemize}
	\begin{figure}
		\centering
		\includegraphics[width=0.6\linewidth]{"Figures/fig 14_11"}
	\end{figure}
\end{frame}

\begin{frame}[wide]{The Troubled Asset Relief Program}
	\begin{itemize}
		\item Economists agree on the importance of restoring the financial system, but debate surrounds the most effective policy
		
		\item Policies such as the Troubled Asset Relief Program (TARP) have been implemented by the Federal Reserve and Department of the Treasury

	\end{itemize}
\end{frame}


\begin{frame}[wide]{The Troubled Asset Relief Program}
	\begin{itemize}
		\item TARP included measures like purchasing toxic assets and injecting capital into financial institutions
		\begin{itemize}
			\item Purchases of ``toxic" assets:	Banks possess bad assets, which limits lending
			\item Capital injections into financial institutions: Original TARP invested \$25 billion in each large financial institution
		\end{itemize}
		\item In the fall of 2008, TARP established a \$700 billion fund for the Treasury to purchase toxic assets
		\begin{itemize}
			\item purchase equity in banks and other financial institutions, to guarantee loans to those institutions, and even to bail out some U.S. automakers
		\end{itemize}
		\item Complete reorganization of financial institutions: reorganize debt into new equity claims for the former debt holders
		\item The program did not actually incur large losses; as of the end of 2012, about 97 percent of the disbursements had been returned to the Treasury 
	\end{itemize}
\end{frame}

\begin{frame}[wide]{Fiscal Stimulus}
	\begin{itemize}
		\item American Recovery and Reinvestment Act (ARRA) 2009
		\begin{itemize}
			\item Tax cuts and new government spending -- More than \$250 billion in tax cuts and More than \$500 billion in new government spending			
			\item Increased the deficit to 10 percent of GDP in 2009 (was 3\% of the GDP in 2008)
		\end{itemize}
		\item Most economists believed there should be a fiscal stimulus
		\item Disagreement stemmed from the type of spending and weight of tax cuts
		\begin{itemize}
			\item The Ricardian equivalence argument was also a factor
			\begin{itemize}
				\item High spending must be financed by higher taxes in the future
				\item The prospect of these taxes may reduce the current impact of the stimulus package
			\end{itemize}
			\item Temporary measures that seek to change the timing of production may overcome these concerns (e.g., investment tax credit, etc)
		\end{itemize}
	\end{itemize}
\end{frame}

\begin{frame}[wide]{Forward Guidance}
	\begin{itemize}
		\item The dashed lines are the market expectations about the future fed funds rate
		\item The market seems to believe the interest rate would stay low for sometime after Jan 2012
	\end{itemize}
	\begin{figure}
		\centering
		\includegraphics[width=0.7\linewidth]{"Figures/fig 14_15"}
	\end{figure}
	
\end{frame}

\begin{frame}[wide]{The European Debt Crisis}
	\begin{itemize}
		\item Financial crisis was not a U.S. phenomenon
		\item Sovereign debt crisis in Europe
		\begin{itemize}
			\item Government debt of countries like Greece, Ireland, Italy, and Spain came under pressure 
			\item Interest rates on the debt rising as high as 7 percent
			\item Severe problems in the banking sector 
			\item High unemployment
			
		\end{itemize}
	\end{itemize}
\end{frame}

\begin{frame}[wide]{Financial Reform}
	\begin{itemize}
		\item A key downside to the way in which the financial crisis was handled by policymakers is that it has made problems of moral hazard more severe
		\begin{itemize}
			\item Many of these institutions were deemed too systemically important to fail or too big to fail
			\item Bankers would expect the banking sector will be bailed and would take any risk
		\end{itemize}
		\item How do we prevent major problems in the financial system?
		\begin{itemize}
			\item Gain greater understanding of volatile prices (e.g., housing, stocks, bubbles)
			\item Acknowledge costs with the benefits of financial intervention
			\item Understand the downside of moral hazard -- minimize incentive for bankers to take ``heads, I win; tails, the economy loses" types of bets
		\end{itemize}
	\end{itemize}
\end{frame}

\begin{frame}[wide]{Financial Reform}
	\begin{itemize}
		\item How do firms fail under normal circumstances?
		\begin{itemize}
			\item A common solution when a firm's liabilities surpass its assets is to declare bankruptcy
			\item A judge or some other government entity steps in and reorganizes the firm if possible so that it can resume its business
			\item Equity (net worth) is already zero or even negative, so the stockholders have lost their entire investment
			\item Debt is written down to zero, and the former debtholders are given equity claims in the newly reorganized firm
			\item The value of the debt at the time of the reorganization becomes the new equity in the new firm
		\end{itemize}
	\end{itemize}
\end{frame}

\begin{frame}[wide]{Financial Reform}
	\begin{itemize}
		\item This approach has a number of appealing features for financial sector as well
		\begin{itemize}
			\item Stockholders and bondholders bear the brunt of the burden, not taxpayers
			\item With the bank recapitalized, it should emerge with the ability to lend
		\end{itemize} 
		\item Potential problem with this approach
		\begin{itemize}
			\item Severely interferes with the functioning of financial markets if not handled carefully
			\item Liquidity problem in the short run
			\item The panic following the collapse of Lehman Brothers in September 2008 is a reminder of what can go wrong
		\end{itemize}
	\end{itemize}
\end{frame}

\begin{frame}[wide]{Financial Reform}
	Squam Lake Group suggested guidelines for financial reform:
	\begin{itemize}
		\item Create a systemic regulator
		\begin{itemize}
			\item Each country should have one primary regulator charged with monitoring systemic risk in financial markets
			\item Responsible for enforcing the other elements of financial reform
		\end{itemize}
		\item Enhance capital requirements
		\begin{itemize}
			\item The minimum percentage of capital requirements (equity/capital to asset ratio) should be raised so that financial institutions have a larger “cushion” in case the value of their assets declines suddenly
			\item Also, larger or more complex financial institutions should face higher capital requirements
		\end{itemize}
		\item Main cost of bailouts
		\begin{itemize}
			\item Borne by people outside of the finance industry
			\item Measured by lost jobs and forgone GDP
		\end{itemize}
	\end{itemize}
\end{frame}

\begin{frame}[wide]{Financial Reform}
	\begin{itemize}
		\item Link executive compensation to long-term performance
		\begin{itemize}
			\item The compensation of senior management in systemically important financial institutions should be tied, at least in part, to the long-term performance of the firms
		\end{itemize}
		\item Require “living wills.” 
		\begin{itemize}
			\item Some important financial institutions will inevitably fail
			\item A living will—a set of instructions for how this failure should be carried out to keep to a low level the disruption to financial markets—would minimize the effects of such a failure
%			\item The direct cost of bailing out financial institutions was only a small portion of the overall cost of the global financial crisis
%			\item The bulk of these costs was borne by people outside the finance industry
		\end{itemize}
	\end{itemize}
\end{frame}

\begin{frame}[wide]{Financial Reform}
	\begin{itemize}
		\item Require convertible debt
		\begin{itemize}
			\item Systemically important financial institutions might also be required to issue a significant amount of a new, hybrid debt instrument that converts to equity when certain prespecified crisis conditions arise
			\begin{itemize}
				\item This hybrid is called convertible debt
			\end{itemize}
			\item When the value of equity in such a firm falls significantly, this debt would convert to equity, recapitalizing banks at no cost to taxpayers
			\item Bondholders and shareholders would themselves bear the costs of financial problems, minimizing the likelihood of taxpayer bailouts
			\item This convertible debt functions like an automated, limited version of the reorganization that occurs in bankruptcy, but it should lesser effect on the financial markets that a more formal reorganization would 
		\end{itemize} 
	\end{itemize}
\end{frame}

\begin{frame}[wide]{The Aftermath of the Great Recession}
	The slow recovery following the Great Recession has prompted two proposed explanations:
	\begin{itemize}
		\item Secular stagnation, characterized by sluggish economic growth, persistently low inflation, and negative real interest rates
		\begin{itemize}
			\item Japan has experienced with this phenomenon since the 1990s
		\end{itemize}
		\item The combination of secular stagnation and a "global saving glut" has depressed investment demand while simultaneously increasing the supply of saving, resulting in negative real interest rates
		\item The presence of the zero lower bound does not allow the market to be cleared
		\item Reduced investment leads to less capital accumulation. As indicated in the Solow model, the economic grow slower
	\end{itemize}
	
\end{frame}

\begin{frame}[wide]{The Aftermath of the Great Recession}
	\begin{itemize}
		\item TFP growth has been slow since 2004
		\item R\& D spending remains strong. But, it may be because the financial crisis increased the misallocation of resources
	\end{itemize}
	\begin{figure}
		\centering
		\includegraphics[width=0.5\linewidth]{"Figures/fig 14_16"}
	\end{figure}
	
\end{frame}

\begin{frame}[wide]{Macroeconomic Research after the Financial Crisis}
	\begin{itemize}
		\item Did macroeconomists fail? 
		\begin{itemize}
			\item Research was not available to provide guidelines
			\item Research has been limited by computational power and data availability
		\end{itemize}
		\item New models are more advanced and include
		\begin{itemize}
			\item Frictions: sticky inflation and limited financial markets
			\item Heterogeneous consumers -- some are rich, some are poor; some are more willing to save than the other, etc
			\item Imperfect financial markets -- agents cannot buy any asset or combination of asset such that it would cancel out all negative shock
			\item Allowing for more roles of government
		\end{itemize}
		\item Economist response with a sharp increase in macroeconomic research focused on economic fluctuations and financial imperfections
	\end{itemize}
\end{frame}

\begin{frame}[wide]{Conclusion}
	\begin{itemize}
		\item The global financial crisis will likely go down in history as the largest recession since the Great Depression
		\item The “Great Recession” was a balance sheet crisis for
		\begin{itemize}
			\item Financial institutions
			\item Households
			\item The Federal Reserve
			\item The government
		\end{itemize}
		\item As we recover from the crisis, the question remain, is the productivity slow down temporary?
		\begin{itemize}
			\item The experience from Great Depression would suggest no...
		\end{itemize}
	\end{itemize}
\end{frame}

%\begin{frame}[wide]{An Increase in the Investment Rate}
%	\begin{columns} 
%			% Column 1
%			\begin{column}{.5\textwidth}
%				\begin{itemize}
%					\item Suppose the investment rate increases permanently for exogenous reasons ($s'>s$)
%					\begin{itemize}
%						\item The investment curve rotates upward ($\bar{s}'\bar{A}  \bar{L}^{1-\alpha} K_t^\alpha > \bar{s}\bar{A}  \bar{L}^{1-\alpha} K_t^\alpha$)
%						\item The depreciation curve remains unchanged ($\delta K_t$)
%						\item The capital stock increases by transition dynamics to reach the new steady state
%						\begin{itemize}
%							\item This happens because investment exceeds depreciation ($sY_t>\delta K_t$)
%						\end{itemize}
%						\item The new steady state is located to the right
%					\end{itemize}
%				\end{itemize}
%			\end{column}
%			% Column 2    
%			\begin{column}{.5\textwidth}			
%				\begin{figure}
%					\centering
%					\includegraphics[width=1\linewidth]{"Figures/An Increase in the Investment Rate"}
%				\end{figure}
%			\end{column}%
%		\end{columns}
%\end{frame}

%\begin{frame}[wide]{Measuring the State of the Economy}
%	\begin{columns} 
%		% Column 1
%		\begin{column}{.5\textwidth}
%			\begin{table}[htbp]
%				\centering
%				\begin{tabular}{lrrrr}
%					& 2005  & 2006  & 2007  & 2008 \\
%					\hline
%					GDP   & 12.6  & 13.4  & 14.0    & 14.3 \\
%					(in trillions of \$) &       &       &       &  \\
%					Growth rate GDP &       & 0.063 & 0.045 & 0.021 \\
%					\hline
%				\end{tabular}%
%				\label{tab:addlabel}%
%			\end{table}%
%		\end{column}
%		% Column 2    
%		\begin{column}{.5\textwidth}			
%			
%			
%		\end{column}%
%	\end{columns}
%\end{frame}

\begin{frame}[wide]{Next Lecture}
	\begin{itemize}
		\item Next lecture will cover Ch.15
	\end{itemize}
\end{frame}

\end{document}